\section{Problem Summary}
Use every number from \(1\) through \(n\) exactly once and split them into two sets with the same total sum. If this cannot be done, print that immediately.

\section{Editorial}
The total sum is
\[
\frac{n(n + 1)}{2}.
\]
If that sum is odd, then two equal halves are impossible.

Assume the total sum is even and let the target be half of it. We can build one set greedily from large numbers down to small numbers:
\begin{itemize}
\item inspect \(n, n - 1, \ldots, 1\)
\item whenever the current number does not exceed the remaining target, take it into the first set
\item subtract it from the remaining target
\end{itemize}

This works because the numbers are consecutive. If a large number fits, taking it is always safe and makes the remaining target smaller; if it does not fit, we simply skip it and keep looking. By the time we reach \(1\), the remaining target must be \(0\).

Everything that was not taken belongs to the second set.

\section{Pseudocode}
Compute the total sum of \(1\) through \(n\).

If the total sum is odd:
\begin{itemize}
\item print \texttt{NO}
\item stop
\end{itemize}

Set the remaining target to half of the total.

For values from \(n\) down to \(1\):
\begin{itemize}
\item if the value is at most the remaining target, put it in the first set and subtract it from the target
\item otherwise leave it for the second set
\end{itemize}

Print both sets.

\section{Complexity Analysis}
\begin{itemize}
\item Time: \(O(n)\)
\item Memory: \(O(n)\) to store the chosen partition
\end{itemize}

\section{Notes}
Any valid partition is accepted, so the greedy descending construction is enough.
